% !TeX program = xelatex
\documentclass[aspectratio=169]{beamer}

\usepackage{import}

\subimport{../../}{system.tex}
\UseTemplateSet{
  layout = beamer,
  globalfonts = {cmu,shanggu},
  fonts = {phoenician,aramaic,pahlavi_parthian,pahlavi_inscriptional,pahlavi_psalter,avestan,kharosthi,samaritan,hebrew},
  features = {tables,image}
}

\title{亞非系文字}
\subtitle{世界文字總覽~第二講}
\author{Kelvin Yang}
\date{\today}

\begin{document}
\begin{frame}[plain]
    \titlepage
\end{frame}
\AgendaFrame

\section{亞非系文字}
\begin{frame}{亞非系文字}
    本講對亞非系文字進行概述，並重點介紹希伯來文、阿拉伯文以及滿文。

    \begin{block}{亞非系文字}
        \textbf{亞非系文字}（Afroaisatic scripts）即包含腓尼基文、古希伯來文、阿拉米文及其相關衍生文字在内的、以輔音音位文字爲主的、通行於西亞--北非到中亞廣大區域的、以亞非語族羣爲主的人羣所使用的一系列文字。
    \end{block}

    總體上，亞非系文字可以分爲三個子系：

    \begin{itemize}
        \item \textbf{希伯來子系}（Hebraic）
        \item \textbf{阿拉伯子系}（Arabic）
        \item \textbf{敘利亞子系}（Syriac）
    \end{itemize}

    除了這三子系之外，還有一些相對零散、獨立，如缽羅婆文、阿維斯陀文、佉盧文等。
\end{frame}

\begin{frame}{腓尼基文與古希伯來文}
    \textbf{腓尼基文}（Phoenician）大約產生於公元前二千紀末的黎凡特地區，用於拼寫亞非語系的腓尼基語。腓尼基人海上貿易發達，對文字的需求相對較大，而當時的其他主要文字書寫起來都並不方便。
    於是，腓尼基人借用了西北閃米特人（Northwest Semitic）從埃及僧侶體（Hieratic）改造而來的原始西奈文（Proto-Sinaitic），創製了腓尼基文。

    \textbf{古希伯來文}（Paleo-Hebrew）和腓尼基文同源，且相當相似，某種程度上可以當作同一種文字的不同地方變體，用於書寫聖經希伯來語（Biblical Hebrew）。古希伯來文的使用與古代以色列王國的形成密切相關，多用於地方王國行政、紀念銘刻與宗教社群記憶。

    \begin{block}{文字樣例（塔布尼特石棺銘文，公元前五世紀）}
        {\raggedleft
        \PH{𐤀𐤍𐤊 𐤕𐤁𐤍𐤕 𐤊𐤄𐤍 𐤏𐤔𐤕𐤓𐤕 𐤌𐤋𐤊 𐤑𐤃𐤍𐤌 𐤁𐤍 𐤀𐤔𐤌𐤍𐤏𐤆𐤓 𐤊𐤄𐤍 𐤏𐤔𐤕𐤓𐤕 𐤌𐤋𐤊 𐤑𐤃𐤍𐤌 𐤔𐤊𐤁 𐤁𐤀𐤓𐤍 𐤆}
        }

        ʾnk tbnt khn ʿštrt mlk ṣdnm bn ʾšmnʿzr khn ʿštrt mlk ṣdnm škb bʾrn z

        \rmfamily\itshape
        我，塔布尼特，阿斯塔蒂之祭司、西頓王，以實慕那撒之子，阿斯塔蒂之祭司、西頓王，今安臥於此石棺之中。
    \end{block}
\end{frame}

\begin{frame}{阿拉米文}
    \textbf{阿拉米文}（Aramaic）產生於公元前十世紀左右，來源於腓尼基文，用於拼寫敘利亞和兩河流域一些小國所使用的阿拉米語。
    其發展可以分爲下列階段：

    \begin{itemize}
        \item 公元前十--九世紀，阿拉米文還是使用於阿拉米諸國的地方性文字。
        \item 公元前九--八世紀，新亞述帝國（Neo-Assyrian Empire）征服阿拉米諸國後，阿拉米文的簡便性使之得到廣泛推廣。
        \item 公元前八--四世紀，阿契美尼德王朝時期發展爲帝國阿拉米文（Imperial Aramaic），被用作波斯帝國官方行政文字之一。
        \item 公元前四世紀之後，波斯帝國的瓦解帶來了阿拉米文的分化，演進出多種地方變體。
    \end{itemize}

    \begin{block}{文字樣例}
        \raggedleft
        \IA{𐡊𐡋 𐡁𐡓𐡍𐡔𐡀 𐡁𐡓𐡉𐡋𐡄 𐡇𐡀𐡓𐡀 𐡅𐡁𐡓𐡁𐡓 𐡂𐡅 𐡀𐡉𐡒𐡓𐡀 𐡅𐡆𐡃𐡒𐡀.}
    \end{block}
\end{frame}

\begin{frame}{阿拉米文}
    \begin{OneImage}[htbp][\linewidth][0.75\textheight]{_assets/20260310_195239.png}[埃利芬汀紙草文書（Elephantine Papyri），前407年]\end{OneImage}
\end{frame}

\begin{frame}{缽羅婆文}
    \textbf{缽羅婆文}（Pahlavi）是一些用於書寫中古伊朗語的文字，使用於公元前二--公元七世紀，由帝國阿拉米文衍生而來，是薩珊波斯帝國（Sasanian Empire）的官方文字。
    缽羅婆文仍然是輔音音位文字，其中有四種比較重要：
    \begin{itemize}
        \item 碑銘安息文（Inscriptional Parthian）：使用於公元二--五世紀的安息帝國（Parthian Empire）與薩刪帝國早期，用於書寫安息語。
        \item 碑銘缽羅婆文（Inscriptional Pahlavi）：使用於公元三--六世紀薩珊波斯帝國的一些碑銘中，用於拼寫中古波斯語（Middle Persian）。
        \item 詩篇缽羅婆文（Psalter Pahlavi）：發現於公元六七世紀左右的敘利亞語詩篇中，需要連寫。
        \item 書卷缽羅婆文（Book Pahlavi）：最常見的缽羅婆文，有複雜的連寫形式，一直使用到約公元900年。
    \end{itemize}
\end{frame}

\begin{frame}{缽羅婆文}
    公元四--六世紀左右，薩珊王朝沙普爾二世（Shapur II）統治時期，爲了準確記錄阿維斯陀語的祆教（Zoroastrian）經典，在缽羅婆文的基礎上創製出了\textbf{阿維斯陀文}（Avestan）。
    阿維斯陀文不再是輔音音位文字，而是一種全音位文字。

    \begin{block}{文字樣例}
        碑銘安息文：\hfill\PAR{𐭊𐭋 𐭁𐭓𐭍𐭔𐭀 𐭁𐭓𐭉𐭋𐭄 𐭇𐭀𐭓𐭀 𐭅𐭁𐭓𐭁𐭓 𐭂𐭅 𐭀𐭉𐭒𐭓𐭀 𐭅𐭆𐭃𐭒𐭀.}
        
        碑銘缽羅婆文：\hfill\PAH{𐭪𐭫 𐭡𐭥𐭭𐭱𐭠 𐭡𐭥𐭩𐭫𐭤 𐭧𐭠𐭥𐭠 𐭥𐭡𐭥𐭡𐭥 𐭢𐭥 𐭠𐭩𐭬𐭥𐭠 𐭥𐭦𐭣𐭬𐭠.}

        詩篇缽羅婆文：\hfill\PSP{𐮉𐮊 𐮁𐮅𐮌𐮐𐮀 𐮁𐮅𐮈𐮊𐮄 𐮇𐮀𐮅𐮀 𐮅𐮁𐮅𐮁𐮅 𐮂𐮅 𐮀𐮈𐮋𐮅𐮀 𐮅𐮆𐮃𐮋𐮀.}

        阿維斯陀文：\hfill\AV{𐬞𐬎𐬭𐬕𐬍𐬝 𐬛𐬁𐬢𐬁 𐬋 𐬨𐬀𐬌𐬢𐬌𐬌𐬋 𐬑𐬀𐬭𐬛 𐬐𐬎 𐬐𐬋𐬭𐬏𐬀𐬱𐬨 𐬬𐬀𐬝𐬙𐬀𐬭 𐬀𐬌𐬌𐬃 𐬐𐬋𐬭𐬛𐬌𐬮}
    \end{block}
\end{frame}

\begin{frame}{佉盧文}
    \textbf{佉盧文}（Kharosthi）是唯一不來自婆羅米文的印度文字，使用於公元前四--公元三世紀之間的犍陀羅（Gandhara）地區，是一種元音附標文字。
    其具體的來源並不清楚，但可能是從阿拉米文演變而來，用於拼寫多種印度語言，主要是犍陀羅語（Gandhari）。

    \begin{block}{文字樣例}
        \raggedleft
        \KH{𐨩𐨟𐨿𐨪 𐨗𐨒𐨟𐨁 𐨭𐨌𐨣𐨿𐨟𐨁𐨣𐨿𐨩𐨌𐨩𐨯𐨿𐨬𐨌𐨟𐨣𐨿𐨟𐨿𐨪𐨿𐨩𐨌𐨞𐨌𐨎 𐨀𐨌𐨢𐨌𐨪𐨏 𐨨𐨌𐨣𐨬𐨤𐨪𐨁𐨬𐨌𐨪𐨯𐨿𐨩 𐨯𐨪𐨿𐨬𐨅𐨮𐨌𐨨𐨤𐨁}
    \end{block}
\end{frame}

\section{希伯來子系}
\begin{frame}{希伯來子系}
    \textbf{希伯來子系}（Hebraic）主要指古希伯來文的衍生文字以及方體阿拉米文即現代希伯來文。

    \textbf{撒馬利亞文}（Samaritan）直接來源於古希伯來文，由撒馬利亞人用於書寫撒馬利亞語，大約創製於公元前四世紀，一直使用至今。

    \begin{block}{文本樣例}
        \raggedleft
        \SM{ࠊࠋ ࠁࠓࠍࠔࠀ ࠁࠓࠉࠋࠄ ࠇࠀࠓࠀ ࠅࠁࠓࠁࠓ ࠂࠅ ࠀࠉࠒࠓࠀ ࠅࠆࠃࠒࠀ. }
    \end{block}

    \textbf{希伯來文}（Hebrew），實際上就是方體阿拉米文（Square Aramaic），是阿拉米文的一種地方變體。
    在巴比倫之囚（Babylonian captivity）後，大量古猶太人被擄去巴比倫，進入了後來波斯帝國的統治範圍。
    於是，在公元前六--二世紀之間，這些猶太人逐漸由原本的古希伯來文轉而使用阿拉米文，到公元前一世紀左右基本定型。
\end{frame}

\begin{frame}{希伯來文}
    我們選取現代俄語來對希伯來文進行介紹。

        \small
    {\setlength{\columnsep}{0.6em}% ← 三欄之間距離：0.4em~0.8em 自己微調
    \setlength{\tabcolsep}{3pt}%  ← 表內欄距：2pt~4pt 自己微調

    \begin{columns}[T]
        \column{0.3\textwidth}
            \centering
            \begin{tabularx}{\linewidth}{@{}rrl@{}}
                \toprule
                \NiceHead{ 字母 & 名稱 & IPA \\ }
                \HE{א} & \HE{אָ֫לֶפ} & /ʔ/ \\
                \bottomrule
            \end{tabularx}

        \column{0.25\textwidth}
            \centering
            \begin{tabularx}{\linewidth}{@{}lllc@{}}
                \toprule
                \NiceHead{ 字母 & 名稱 & IPA & 對照 \\ }

                \bottomrule
            \end{tabularx}

        \column{0.4\textwidth}
            \centering
            \begin{tabularx}{\linewidth}{@{}lllc@{}}
                \toprule
                \NiceHead{ 字母 & 名稱 & IPA & 對照\\ }
                \bottomrule
            \end{tabularx}
    \end{columns}
    }
\end{frame}

\end{document}
